\chapter{Introduction to calculating with Physical Quantities}

\begin{definicion}[Quantity]
    A quantity is a property of a phenomenon, body or substance that may be distinguished qualitatively and determined quantitatively. It is expressed as the \emph{product of a numerical value and a unit.}
\end{definicion}

There are two types of quantities:
\begin{description}
    \item[Scalar] A quantity that is completely described by its numerical value and unit. Examples include mass, time, temperature, etc.
    \item[Vectorial] A quantity that is completely described by its numerical value (module), unit and direction. Examples include velocity, force, etc.
    
    Usually, only the module of a vectorial quantity is given, and the direction is determined by the context. Other times, the direction is given by a unit vector, which is a vector of module 1 whose only purpose is to indicate the direction of the quantity.
\end{description}

In order to work with quantities, we need to introduce the concept of coordinate system.

\section{Coordinate Systems}

A coordinate system of $\bb{R}^n$ is defined by a point $O$ called the origin and a base $\cc{B}=\{\vec{e}_1, \vec{e}_2, \ldots, \vec{e}_n\}$ of $\bb{R}^n$. The base vectors $\vec{e}_i$ are called the base vectors of the coordinate system. The coordinates of a point $P$ in the coordinate system are the coefficients of the linear combination of the base vectors that gives the vector $\vec{OP}$:
\begin{equation*}
    \vec{OP} = \sum_{i=1}^n x_i \vec{e}_i
\end{equation*}

The basic coordinate system is the cartesian coordinate system, which is defined by the origin $O=(0,0,0)$ and the base $\cc{B}_u=\{\hat{i}, \hat{j}, \hat{k}\}$, where $\hat{i}$, $\hat{j}$ and $\hat{k}$ are the unit vectors in the direction of the $x$, $y$ and $z$ axes, respectively. They should always follow the right-hand rule:
\begin{equation*}
    \hat{i} \times \hat{j} = \hat{k}
\end{equation*}

However, there are other coordinate systems that are more convenient to work with in certain situations. In those cases, they are usually not given by a base, but by a transformation that relates the coordinates of the new system with the coordinates of the cartesian system. In those cases, each vector of the new system is given by fixing all the coordinates of the new system except one, and letting that coordinate vary. More formally, it the new system is given by the coordinates ($u_1, u_2, u_3$), then the vector $\vec{e}_i$ of the new system is given by:
\begin{equation*}
    \vec{e}_i = \frac{\partial \vec{r}}{\partial u_i}
\end{equation*}
where $\vec{r}$ is the position vector of a point.

\subsection{Polar coordinates}

The polar coordinate system (two dimensions) is given by the coordinates $(\rho, \phi)$, where $\rho$ is the distance from the origin and $\phi$ is the angle between the positive $x$ axis and the line that connects the origin with the point. It is shown in Figure~\ref{fig:polar_coordinates}. The transformation that relates the polar coordinates with the cartesian coordinates is given by:
\begin{equation*}
    \begin{cases}
        x = \rho \cos \phi \\
        y = \rho \sin \phi
    \end{cases}
\end{equation*}
\begin{figure}
    \centering
    \begin{tikzpicture}

        % --- Definición de Parámetros (puedes cambiarlos) ---
        % Longitud del radio (r) y ángulo en grados (t)
        \def\r{3.5} % Longitud del radio
        \def\theta{40} % Ángulo en grados

        % --- Ejes Cartesianos ---
        \draw[-Stealth, thick, gray!60!black] (-0.5,0) -- (4.5,0) node[right, black] {$x$};
        \draw[-Stealth, thick, gray!60!black] (0,-0.5) -- (0,4) node[above, black] {$y$};
        \node[anchor=north east] (O) at (0,0) {$O$}; % Origen

        % --- Punto P y Líneas Polares ---
        % Definir el punto P(r, theta) en coordenadas polares
        \coordinate (P) at (\theta:\r); 

        % Dibujar la línea punteada (el radio r) desde el origen hasta P
        \draw[dashed, blue!70!black, thick] (0,0) -- (P) node[midway, sloped, above, font=\small] {$r$};

        % Dibujar el arco que representa el ángulo theta
        % (Usamos la librería 'angles' para esto automáticamente)
        \coordinate (Xaxis) at (4,0); % Punto de referencia en el eje 
        \coordinate (Oaux) at (0,0);

        \draw pic[
                -Stealth,
                draw=blue!70!black,
                angle radius=1.2cm,
                angle eccentricity=1.3,
                pic text={\color{blue!70!black}$\phi$}
            ] {angle = Xaxis--Oaux--P};

        % Dibujar el punto P y su etiqueta
        \filldraw[blue!70!black] (P) circle (2pt);
        \node[blue!70!black, anchor=west, font=\boldmath] at (P) {};

        % --- Proyecciones Cartesianas ---
        % Proyección sobre el eje X (x = r cos(t))
        \draw[dotted, gray] (P) -- ({ \r*cos(\theta) }, 0);
        \node[anchor=north, font=\small] at ({ \r*cos(\theta) }, 0) {$x$};

        % Proyección sobre el eje Y (y = r sin(t))
        \draw[dotted, gray] (P) -- (0, { \r*sin(\theta) });
        \node[anchor=east, font=\small] at (0, { \r*sin(\theta) }) {$y$};

        % Cuarto de circulo de radio r para mostrar la dirección del ángulo
        \draw[dotted] (\r,0) arc (0:90:\r);

        % Vectores unitarios en coordenadas polares
        \draw[-Stealth, thick, red] (P) -- ++({cos(\theta)}, {sin(\theta)}) node[above, red] {$\vec{e}_\rho$};
        \draw[-Stealth, thick, red] (P) -- ++({-sin(\theta)}, {cos(\theta)}) node[above, red] {$\vec{e}_\phi$};

    \end{tikzpicture}
    \caption{Polar coordinates.}
    \label{fig:polar_coordinates}
\end{figure}


In order to calculate the base vectors of the polar coordinate system, we need to calculate the position vector of a point in polar coordinates. The position vector of a point $P$ in polar coordinates is given by:
\begin{equation*}
    \vec{r} = x \hat{i} + y \hat{j}
    = \rho \cos \phi\  \hat{i} + \rho \sin \phi\  \hat{j}
\end{equation*}

Therefore:
\begin{align*}
    \vec{e}_\rho &= \dfrac{\cos \phi\  \hat{i} + \sin \phi\  \hat{j}}{\sqrt{\cos^2 \phi + \sin^2 \phi}} = \cos \phi\  \hat{i} + \sin \phi\  \hat{j} \\
    \vec{e}_\phi &= \dfrac{-\rho \sin \phi\  \hat{i} + \rho \cos \phi\  \hat{j}}{\sqrt{(-\rho \sin \phi)^2 + (\rho \cos \phi)^2}} = -\sin \phi\  \hat{i} + \cos \phi\  \hat{j}
\end{align*}

As we can see, the base vectors of the polar coordinate system are not constant, but they depend on the angle $\phi$. This is because the polar coordinate system is a curvilinear coordinate system, where the axes are not straight lines, but curves. Therefore, the base vectors of a curvilinear coordinate system are not constant, but they depend on the coordinates of the point.

\subsection{Cylindrical coordinates}

The cylindrical coordinate system (three dimensions) is given by the coordinates $(\rho, \phi, z)$, where $\rho$ and $\phi$ are the same as in the polar coordinate system when projected onto the $xy$ plane, and $z$ is the height above the $xy$ plane.
They are shown in Figure~\ref{fig:cylindrical_coordinates}.
The transformation that relates the cylindrical coordinates with the cartesian coordinates is given by:
\begin{equation*}
    \begin{cases}
        x = \rho \cos \phi \\
        y = \rho \sin \phi \\
        z = z
    \end{cases}
\end{equation*}
\begin{figure}
    \centering
    \tdplotsetmaincoords{70}{110}
\begin{tikzpicture}

    % --- Configuración de Parámetros Modificables ---
    \tikzmath{
        \rVal = 8;          % Reducido un poco para que quepa bien en la vista
        \anglePhi = 35;      
        \angleTheta = 35;    
        \rCil = \rVal * cos(\angleTheta); % Radio del cilindro (proyección en XY)
        \hCil = \rVal * sin(\angleTheta); % Altura del punto P
    }

    \begin{scope}[tdplot_main_coords] 
        % Ejes Cartesianos
        \draw[-Stealth, thick] (0,0,0) -- (7.5,0,0) node[anchor=north east] {$x$};
        \draw[-Stealth, thick] (0,0,0) -- (0,7.5,0) node[anchor=north west] {$y$};
        \draw[-Stealth, thick] (0,0,0) -- (0,0,4) node[anchor=south] {$z$};

        % --- Puntos Clave ---
        \coordinate (O) at (0,0,0);
        \coordinate (P) at ({\rCil*cos(\anglePhi)}, {\rCil*sin(\anglePhi)}, \hCil);
        \coordinate (Pxy) at ({\rCil*cos(\anglePhi)}, {\rCil*sin(\anglePhi)}, 0);
        \coordinate (Xorigin) at (1.5,0,0);

        % --- Cilindro y Proyección ---
        
        % 1. Proyección del círculo en el plano XY (el arco de 0 a 90 grados)
        \draw[blue!20!black, thick, opacity=0.5] (\rCil,0,0) arc (0:90:\rCil);
        
        % 2. "Pared" del cilindro (arco a la altura de P)
        \draw[blue!20!black, thick, opacity=0.5] (\rCil,0,\hCil) arc (0:90:\rCil);
        
        % 3. Líneas verticales del cilindro (bordes)
        \draw[blue!20!black, opacity=0.3] (\rCil,0,0) -- (\rCil,0,\hCil+1);
        \draw[blue!20!black, opacity=0.3] (0,\rCil,0) -- (0,\rCil,\hCil+1);

        % --- Líneas de Referencia del Punto ---
        \draw[thick, black!90] (O) -- (P) node[midway, sloped, above] {$r$};
        \draw[dashed, black!60] (Pxy) -- node[midway, right] {$z$} (P);
        \draw[dashed, black!60] (O) -- (Pxy) node[midway, sloped, below, font=\small] {$\rho$};

        % --- Ángulos ---
        \draw pic[
            -Stealth, draw=red!70!black, angle radius=0.9cm,
            angle eccentricity=1.3, pic text={\color{red!70!black}$\phi$}
        ] {angle = Xorigin--O--Pxy};

        % --- Base Local en P ---
        \def\baseLength{1.2} 
        \draw[-Stealth, blue!80!black, very thick] (P) -- ++(0,0,\baseLength) node[anchor=south] {$\vec{e}_z$};
        \draw[-Stealth, blue!80!black, very thick] (P) -- ++({\baseLength*cos(\anglePhi)}, {\baseLength*sin(\anglePhi)}, 0) node[anchor=north west] {$\vec{e}_{\rho}$};
        \draw[-Stealth, blue!80!black, very thick] (P) -- ++({\baseLength*cos(\anglePhi+90)}, {\baseLength*sin(\anglePhi+90)}, 0) node[anchor=west] {$\vec{e}_{\phi}$};

        % Puntos
        \filldraw (P) circle (1.5pt);
        \filldraw (Pxy) circle (1.5pt);
        
    \end{scope}
\end{tikzpicture}
    \caption{Cylindrical coordinates.}
    \label{fig:cylindrical_coordinates}
\end{figure}

The base vectors of the cylindrical coordinate system are given by:
\begin{align*}
    \vec{e}_\rho &= \cos \phi\  \hat{i} + \sin \phi\  \hat{j} \\
    \vec{e}_\phi &= -\sin \phi\  \hat{i} + \cos \phi\  \hat{j} \\
    \vec{e}_z &= \hat{k}
\end{align*}

\subsection{Spherical coordinates}

The spherical coordinate system (three dimensions) is given by the coordinates $(r, \theta, \phi)$, where $r$ is the distance from the origin, $\theta$ is the angle between the positive $z$ axis and the line that connects the origin with the point, and $\phi$ is the angle between the positive $x$ axis and the projection of the line that connects the origin with the point onto the $xy$ plane. It is shown in Figure~\ref{fig:spherical_coordinates}.
The transformation that relates the spherical coordinates with the cartesian coordinates is given by:
\begin{equation*}
    \begin{cases}
        x = r \sin \theta \cos \phi \\
        y = r \sin \theta \sin \phi \\
        z = r \cos \theta
    \end{cases}
\end{equation*}
\begin{figure}
    \centering
    \tdplotsetmaincoords{70}{110}
\begin{tikzpicture}

    % --- Parámetros ---
    \def\rVal{5}           
    \def\angPhi{40}        
    \def\angTheta{45}      

    \begin{scope}[x={(-0.35cm,-0.35cm)}, y={(1cm,0cm)}, z={(0cm,1cm)}]
        
        % Ejes Cartesianos
        \draw[-Stealth, thick] (0,0,0) -- (6.5,0,0) node[anchor=north east] {$x$};
        \draw[-Stealth, thick] (0,0,0) -- (0,6.5,0) node[anchor=north west] {$y$};
        \draw[-Stealth, thick] (0,0,0) -- (0,0,5.5) node[anchor=south] {$z$};

        % --- Puntos Clave ---
        \coordinate (O) at (0,0,0);
        \coordinate (P) at ({\rVal*sin(\angTheta)*cos(\angPhi)}, {\rVal*sin(\angTheta)*sin(\angPhi)}, {\rVal*cos(\angTheta)});
        \coordinate (Pxy) at ({\rVal*sin(\angTheta)*cos(\angPhi)}, {\rVal*sin(\angTheta)*sin(\angPhi)}, 0);

        % --- Estructura de la Esfera ---
        
        % 1. Ecuador (Plano XY) - De 0 a 90 grados
        \draw[blue!20!black, opacity=0.3] (\rVal,0,0) arc (0:90:\rVal);
        
        % 2. Meridiano de P (Arco vertical que pasa por P)
        % Para dibujarlo manualmente: rotamos en el plano vertical definido por angPhi
        \draw[blue!20!black, opacity=0.4, thick] 
            plot[domain=0:90, samples=30] 
            ({\rVal*sin(\x)*cos(\angPhi)}, {\rVal*sin(\x)*sin(\angPhi)}, {\rVal*cos(\x)});

        % 3. Paralelo de P (Círculo horizontal)
        \def\rLat{\rVal*sin(\angTheta)}
        \def\zLat{\rVal*cos(\angTheta)}
        \draw[blue!20!black, opacity=0.4, thick] ({\rLat},0,{\zLat}) arc (0:90:{\rLat});

        % --- Líneas de Referencia ---
        \draw[thick] (O) -- (P) node[midway, sloped, above right] {$r$};
        \draw[dashed, gray] (P) -- (Pxy);
        \draw[dashed, gray] (O) -- (Pxy);

        % --- Ángulos ---
        \coordinate (Xaxis) at (1.5,0,0);
        \draw pic[
            -Stealth, draw=red!70!black, angle radius=1cm,
            angle eccentricity=1.3, pic text={\color{red!70!black}$\phi$}
        ] {angle = Xaxis--O--Pxy};

        \coordinate (Zaxis) at (0,0,1.5);
        \draw pic[
            -Stealth, draw=blue!70!black, angle radius=1.0cm,
            angle eccentricity=1.3, pic text={\color{blue!70!black}$\theta$}
        ] {angle = P--O--Zaxis};

        % --- Base Local ---
        \def\bL{1.5}
        % e_r
        \draw[-Stealth, orange!80!black, very thick] (P) -- ++({\bL*sin(\angTheta)*cos(\angPhi)}, {\bL*sin(\angTheta)*sin(\angPhi)}, {\bL*cos(\angTheta)}) node[anchor=south west] {$\vec{e}_r$};
        % e_theta
        \draw[-Stealth, orange!80!black, very thick] (P) -- ++({\bL*cos(\angTheta)*cos(\angPhi)}, {\bL*cos(\angTheta)*sin(\angPhi)}, {-\bL*sin(\angTheta)}) node[anchor=north] {$\vec{e}_\theta$};
        % e_phi
        \draw[-Stealth, orange!80!black, very thick] (P) -- ++({\bL*-sin(\angPhi)}, {\bL*cos(\angPhi)}, 0) node[anchor=west] {$\vec{e}_\phi$};

        \filldraw (P) circle (1.5pt);
        \filldraw[gray] (Pxy) circle (1pt);
        
    \end{scope}
\end{tikzpicture}
    \caption{Spherical coordinates.}
    \label{fig:spherical_coordinates}
\end{figure}

The base vectors of the spherical coordinate system are given by:
\begin{align*}
    \vec{e}_r &= \sin \theta \cos \phi\ \hat{i} + \sin \theta \sin \phi\  \hat{j} + \cos \theta\  \hat{k} \\
    \vec{e}_\theta &= \cos \theta \cos \phi\ \hat{i} + \cos \theta \sin \phi\  \hat{j} - \sin \theta\  \hat{k} \\
    \vec{e}_\phi &= -\sin \phi\  \hat{i} + \cos \phi\  \hat{j}
\end{align*}

